\structuralpart{ВВЕДЕНИЕ}

Рекурсивным называется алгоритм, который в процессе выполнения обращается
к самому себе непосредственно или через другие алгоритмы. Нерекурсивный
алгоритм обходится без таких обращений и повторяет действия с помощью циклов.
Графовой моделью алгоритма называется ориентированный граф, вершины которого
соответствуют операторам алгоритма или их однократным выполнениям,
а дуги~--- связям между ними по передаче управления или данных.

Цель работы~--- провести сравнение рекурсивного и нерекурсивного алгоритмов
на основе графовых моделей алгоритмов.

Для достижения цели необходимо решить следующие задачи:
\begin{itemize}
    \item дать определения четырёх графовых моделей алгоритмов,
    операционного и информационного отношений;
    \item разработать нерекурсивный и рекурсивный последовательные алгоритмы
    решения задачи варианта;
    \item реализовать разработанные алгоритмы;
    \item провести тестирование реализаций алгоритмов;
    \item описать каждый алгоритм четырьмя графовыми моделями;
    \item обосновать, какие фрагменты алгоритмов или их реализаций могут быть
    выполнены в параллельном режиме, а какие не могут.
\end{itemize}

Задача варианта~--- определение количества $k$-значных натуральных чисел
с заданной суммой цифр $s$.
